%! AP Statistics — Unit 3, Lesson 3.2 — Warm-Up
\documentclass[10pt]{article}
\usepackage{apstats-article}
\usepackage{apstats-boxes}

\begin{document}

\pageheader{Unit 3, Lesson 3.2}{Warm-Up}
\namedateperiod

\vspace{0.3in}

% ── SECTION 1: Spiral Review ──────────────────────────────────────────────────
{\large\bfseries\color{navy} Part 1 \quad Spiral Review}

\vspace{0.12in}

\begin{enumerate}

    \item From Lesson 3.1 --- for each data collection scenario, circle whether chance
    was used and whether you would trust the conclusion for the whole population.

    \vspace{8pt}
    \renewcommand{\arraystretch}{1.8}
    \begin{tabularx}{\linewidth}{@{} X p{2.2cm} p{2.8cm} @{}}
    \textbf{Scenario} & \textbf{Chance?} & \textbf{Trustworthy?} \\
    \hline
    A school randomly selects 50 students from the enrollment list and surveys them. &
    Yes \quad No & Yes \quad No \\
    A researcher asks all of her friends whether they prefer cats or dogs. &
    Yes \quad No & Yes \quad No \\
    A website invites all visitors to vote in a poll about sports. &
    Yes \quad No & Yes \quad No \\
    \end{tabularx}

    \vspace{0.18in}

    \item A study finds that students who eat breakfast have higher GPAs. A reporter
    writes: \textit{``Eating breakfast causes students to earn higher grades.''} Is
    the reporter's conclusion justified? Explain.\\[6pt]
    \writeline\\[2pt]\writeline

\end{enumerate}

\vspace{0.35in}

% ── SECTION 2: Looking Ahead ──────────────────────────────────────────────────
{\large\bfseries\color{navy} Part 2 \quad Looking Ahead}

\vspace{0.12in}

\noindent Today we distinguish \textbf{observational studies} from \textbf{experiments}.

\vspace{0.14in}

\begin{tcolorbox}[colback=greenbg, colframe=greenacc, boxrule=0.8pt, arc=2mm,
    left=4mm, right=4mm, top=3mm, bottom=3mm]
    \small
    \begin{itemize}[leftmargin=1.3em, itemsep=8pt]
  \item \textbf{Observational study}: researchers observe without imposing any treatment.
        Cannot establish causation.
  \item \textbf{Experiment}: researchers randomly assign treatments. Can establish
        causation with random assignment.
  \item A sample can only be generalized to the population from which it was
        \emph{randomly} selected.
    \end{itemize}
\end{tcolorbox}

\vspace{0.2in}

\noindent\textcolor{slate}{\small Write one question you have about observational studies
or experiments:}\\[8pt]
\writeline\\[2pt]\writeline

\end{document}
